\unitopen{14}
\begin{candos}
\item użyć celownika (\de{Dativ}) po czasownikach \de{helfen, gefallen, gehören, danken, schenken};
\item zastąpić rzeczownik zaimkiem w bierniku i celowniku (\de{mich/mir, dich/dir, ihn/ihm, \dots});
\item ułożyć poprawnie kolejność dopełnień (\de{Ich gebe es dir});
\item złożyć życzenia i opowiedzieć o prezencie lub imprezie.
\end{candos}

\section{Słownictwo}
\vocabtable{14}

\section{Zaimki osobowe: trzy przypadki}
\begin{center}\small
\begin{tabular}{@{}llll@{}}
\toprule
\textsf{\textbf{Mianownik}} & \textsf{\textbf{Biernik}} & \textsf{\textbf{Celownik}} & \textsf{\textbf{Znaczenie (celownik)}}\\
\midrule
\de{ich} & \de{mich} & \de{mir} & mi, mnie\\
\de{du} & \de{dich} & \de{dir} & ci, tobie\\
\de{er} & \de{ihn} & \de{ihm} & mu, jemu\\
\de{sie} (ona) & \de{sie} & \de{ihr} & jej\\
\de{es} & \de{es} & \de{ihm} & mu\\
\de{wir} & \de{uns} & \de{uns} & nam\\
\de{ihr} & \de{euch} & \de{euch} & wam\\
\de{sie} (oni) & \de{sie} & \de{ihnen} & im\\
\de{Sie} & \de{Sie} & \de{Ihnen} & Panu/Pani\\
\bottomrule
\end{tabular}
\end{center}
Rodzajniki w celowniku znasz z jednostki 10: \artm{dem}, \artf{der}, \artn{dem}, \artf{den}\,+\,\de{-n} w mnogiej (\de{den Kindern}).

\section{Czasowniki z celownikiem}
\de{helfen, gefallen, gehören, danken, gratulieren, antworten, schmecken, passen, wehtun} łączą się \emph{tylko} z celownikiem:
\de{Ich helfe \textbf{dir}.} \de{Das Buch gehört \textbf{mir}.} \de{Die Jacke gefällt \textbf{ihm}.} \de{Mir tut der Kopf weh.}
Czasowniki z dwoma dopełnieniami (\de{geben, schenken, zeigen, schicken, bringen, erklären}) mają celownik (osoba)
i biernik (rzecz): \de{Ich schenke \textbf{meiner Mutter einen Pullover}.}

\begin{zasada}[Kolejność dopełnień]
Dwa rzeczowniki: \emph{celownik przed biernikiem} (\de{Ich gebe dem Kunden den Vertrag}). Dwa zaimki: \emph{biernik przed
celownikiem} (\de{Ich gebe ihn ihm}). Zaimek przed rzeczownikiem: \de{Ich gebe ihm den Vertrag.} \de{Ich gebe ihn dem Kunden.}
\end{zasada}

\begin{falle}
Po polsku „pomagam tobie” i „pomagam cię” różnią się wyraźnie, a po niemiecku błąd jest typowy: \de{Ich helfe dich}.
Poprawnie: \de{Ich helfe dir}. Podobnie \de{Ich danke dich} (błąd) vs \de{Ich danke dir}. Jeśli czasownik jest w polskim
z celownikiem (\emph{pomagać komu}, \emph{dziękować komu}, \emph{gratulować komu}), w niemieckim zwykle też, ale nie zawsze:
\de{fragen} (kogo?) bierze biernik: \de{Ich frage dich}.
\end{falle}

\begin{szwajcaria}
Wśród Szwajcarów używa się \de{Merci vielmal} i \de{Danke vielmal} dla „bardzo dziękuję”; odpowiedź \de{Gern
geschehen} lub \de{Bitte} to norma, a w dialekcie \de{Gärn gscheh}. Życzenia urodzinowe: \de{Alles Gute zum
Geburtstag!} Prezenty w pracy wręcza się raczej skromne, a zaproszenie do domu przychodzi dopiero po
dłuższej znajomości.
\end{szwajcaria}

\begin{praca}
\de{Kannst du mir den Link schicken?} \de{Ich schicke ihn dir gleich.} \de{Der Chef hat mir gratuliert.}
\de{Das Projekt gefällt dem Kunden.} \de{Wem gehört dieser Laptop?} (komu należy? \de{wem}).
\end{praca}

\section{Dialogi}
\textbf{Dialog 1 --- urodziny koleżanki}
\begin{dialog}
\turn{A}{Hallo Anna! Herzlichen Glückwunsch zum Geburtstag!}{Cześć Anno! Serdeczne gratulacje z okazji urodzin!}
\turn{B}{Danke! Das ist aber nett.}{Dziękuję! To bardzo miłe.}
\turn{A}{Ich habe ein Geschenk für dich. Es ist ein Buch.}{Mam dla ciebie prezent. To jest książka.}
\turn{B}{Das gefällt mir sehr! Danke dir!}{Bardzo mi się podoba! Dziękuję ci!}
\turn{A}{Gern geschehen. Ich wünsche dir einen schönen Tag!}{Nie ma za co. Życzę ci miłego dnia!}
\end{dialog}

\section{Ćwiczenia}
\begin{cwiczenia}
\exitem{Uzupełnij zaimkiem (ja): \de{Kannst du \dots{} helfen?}}{\de{mir}}
\exitem{Uzupełnij zaimkiem (on): \de{Das Buch gehört \dots{}.}}{\de{ihm}}
\exitem{Uzupełnij zaimkiem (ona): \de{Ich schenke \dots{} Blumen.}}{\de{ihr}}
\exitem{Uzupełnij zaimkiem (oni): \de{Ich danke \dots{}.}}{\de{ihnen}}
\exitem{Uzupełnij zaimkiem (biernik, ja): \de{Er ruft \dots{} an.}}{\de{mich}}
\exitem{Uzupełnij rodzajnikiem: \de{Ich gebe \dots{} Kunden den Vertrag.} (der Kunde)}{\de{dem}}
\exitem{Zamień oba na zaimki: \de{Ich gebe dem Kunden den Vertrag.}}{\de{Ich gebe ihn ihm.}}
\exitem{Popraw: \de{Ich helfe dich.}}{\de{Ich helfe dir.}}
\exitem{Przetłumacz: „Podoba mi się ten laptop.”}{\de{Der Laptop gefällt mir.}}
\exitem{Przetłumacz: „Boli mnie głowa.”}{\de{Mir tut der Kopf weh.}}
\exitem{Złóż życzenia urodzinowe koledze.}{\de{Alles Gute zum Geburtstag!} lub \de{Herzlichen Glückwunsch zum Geburtstag!}}
\exitem{Zapytaj, do kogo należy ten laptop.}{\de{Wem gehört dieser Laptop?}}
\end{cwiczenia}

\weekend{Opisz po niemiecku ostatni prezent, który komuś dałeś lub dostałeś (komu, co, dlaczego). Użyj co najmniej
trzech zaimków w celowniku i dwóch w bierniku. Poproś AI o sprawdzenie samych przypadków.}
